\scanpage{30}

% L1–2
De a) on déduit, par récurrence transfinie, pour tt $E\in\operatorname{Ob}\Delta$
\[
 \underline{\mathrm{Hom}}_{\deleted{\unsure{\Delta}}d_i\deleted{\unsure{\rho}}}(C_i,E)\xrightarrow{\sim}
 \underline{\mathrm{Hom}}_d(C,E)\qquad(F\mapsto F\circ\alpha_i)\tag*{(1)}
\]

% L3–4
De b) on déduit de $\overline{\mathrm m}$
\[
 \alpha_i(d)\subset d_i\unsure{\subset|=}\alpha_i(d)\cup
 \bigcup_{\substack{\underline{j}<i\\(S,\lambda)\in\unsure{\rho|\sigma}}}
 \bigcup_{u\in\underline{\mathrm{Hom}}(S,C_j)}
 \tau_{ij}u(\lambda)\qquad\substack{\text{si }i\text{ ordinal}\\\text{limite}}\tag*{(2)}
\]

% L5–6
Prouvons que pour $\omega$ \inserted{ordinal limite} assez grand ($\omega$ un ``petit'' ordinal), $(C_\omega,d_\omega)$ est une solution du Pb. \deleted{Cela signifie} \inserted{Il suffit} qu'on ait
% L7–9
$1^\circ)$ $C_\omega\in\operatorname{Ob}\Delta$ i.e. \quad $a_1)$ $\forall(S,\lambda)\in\sigma$ et $u:S-C_\omega$, existence des lim correspondant aux cônes $\in u(\lambda)$
\par $a_2)$ $\forall(\unsure{R_0},R_1,\ldots\varphi)\in\rho$ \deleted{et $u_1:R_1\to C_\omega$ \inserted{$\delta_1$-exact}, $\exists\,u_2:R_2\to C_\omega$ \inserted{$\delta_2$-exact} et $u_1\varphi_1\simeq u_2\varphi_2$}
\[
 \begin{aligned}
 &\underline{\mathrm{Hom}}_{\delta_1}(R_1,C_\omega)\\[-2pt]
 &\hspace{2em}\unsure{\wr|\simeq}\\[-2pt]
 &\underline{\mathrm{Hom}}_{\delta_2}(R_2,C_\omega)
 \end{aligned}
\]

% L10
$2^\circ)$ $\alpha_{\unsure{i|\omega}}:C\longrightarrow C_\omega$ est $d$ exact

% L11
$3^\circ)$ \deleted{\unsure{$E$} \illegible{1}}
\[
 d_\omega=\alpha_\omega(d)\cup\bigcup_{(S,\lambda)\in\sigma}
 \bigcup_{u\in\underline{\mathrm{Hom}}_\lambda(S,C_\omega)}u(\lambda)
\]

% L12–13
[Car $3^\circ$ implique que pour tte $E\in\operatorname{Ob}\Delta$ on a
\[
 \underline{\mathrm{Hom}}_{\Delta,\alpha_\omega(d)}(C_\omega,E)\deleted{\unsure{\simeq}}=
 \underline{\mathrm{Hom}}_{\deleted{\unsure{\Delta,\alpha_\omega(d)}}d_\omega}(C_\omega,E)
 \xrightarrow[(1)]{\sim}\underline{\mathrm{Hom}}_d(C,E)
\]

% L14, interline
\deleted{\unsure{donc|Alors}} Soit $T_\omega$ l'ensemble des ordinaux $j<\omega$

% L15–16
Pour vérifier $2^\circ)$ il suffit de vérifier l'\unsure{assert.|assertion} plus forte
\[
 \forall i<\omega,\quad\tau_{\omega,i}:C_i-C_\omega\text{ est }d_i\text{-exact}\tag*{(3)}
\]

% L17–23
Ceci résulte à son tour de c) Si $T_\omega$ est grand devant les cardinaux : $\operatorname{card}\operatorname{Fl} I$, où $I$ est une cat. d'indices intervenant dans un des cônes $\in d_i$ ; Or en vertu de (2), ces $I$ \deleted{sont \unsure{toutes}} sont \unsure{ou bien|au plus|de plus} ceux qui interviennent dans $d$, soit dans un $\lambda$ associé à un $(S,\lambda)\in\sigma$ -- c'est un petit ens de petites catégories, donc les $\operatorname{card}\operatorname{Fl}I$ sont majorés par un cardinal $\pi$ -- et il suffit $T_\omega$ grand devant $\pi$.

% L18–19, left (slanted)
% slanted
\margin{\unsure{dégager|désigner} lemme}
% archive sheet number: 29
